\section{导读：这期为什么是一堂自动驾驶技术路线课}

本节先定位这期。郎咸朋 2013--2018 年在百度参与自动驾驶和高精地图，2018 年加入理想汽车，过去十年都在中国自动驾驶一线。这期访谈的价值，不是讲某一家公司的八卦，而是用亲历者视角把自动驾驶的技术路线从“高精地图+激光雷达+规则”讲到“BEV+Transformer”“端到端”“VLM+世界模型+RL”。

\begin{importantbox}{本期核心命题}
自动驾驶十年的主线，是从把车当作“有轨电车”研发，转向把驾驶当作可学习的能力。早期路线依赖高精地图、昂贵激光雷达和规则工程；Tesla 的关键贡献是用视觉、BEV、芯片和车队数据升维解决问题；端到端进一步把人工规则压缩进模型和数据，开始从软件功能走向能力系统。
\end{importantbox}

\begin{knowledgebox}{视觉策略说明}
本视频是固定访谈画面，没有投屏和白板。正文只使用封面作为来源识别；正文图像全部为自制概念图，用来解释自动驾驶技术栈演进、传感器取舍、端到端、L3/L4 责任边界、理想数据闭环和组织压力。
\end{knowledgebox}

\begin{figure}[H]
\centering
\includegraphics[width=0.96\textwidth]{av-ten-year-map.png}
\caption{自动驾驶十年路线图：从高精地图和激光雷达到 BEV/Transformer、端到端和世界模型。自制概念图，依据 00:01:32--00:41:14 对谈内容整理。}
\end{figure}

\begin{knowledgebox}{读图：从工程确定性到模型能力}
图中从 HD Map、LiDAR Stack 到 BEV、End-to-End、World Model。读这张图时要看范式变化：早期用地图和传感器把世界变简单；后来用数据和模型直接学习复杂世界；再往后，系统需要预测未来和闭环优化。
\end{knowledgebox}

\subsection{阅读路线}

本节给出读法。全片可以围绕五个问题展开：第一，为什么高精地图和重激光雷达路线像“有轨电车”？第二，Tesla 为什么坚持视觉和自研芯片？第三，BEV 和 Transformer 解决了什么空间建模问题？第四，端到端为什么意味着“以前做自动驾驶都做错了”？第五，L3、L4、世界模型、RL 和理想的数据/组织闭环如何共同决定下一阶段。

\noindent\begin{tabularx}{\textwidth}{p{0.19\textwidth}p{0.34\textwidth}X}
\toprule
阅读问题 & 访谈中的材料 & 要形成的判断 \\
\midrule
早期路线错在哪里 & 高精地图、激光雷达、规则、ODD 穷举 & 用轨道化思维处理开放道路，难以 scale。 \\
Tesla 为什么重要 & 纯视觉、BEV、芯片、升维解法 & 把自动驾驶从传感器堆料转向数据和模型路线。 \\
端到端改变什么 & 场景无法穷举、规则互相干扰、模型学习能力 & 从功能开发转向能力训练。 \\
L3/L4 怎么理解 & 系统责任、接管、安全区域、ODD & L3 不是 L2 延长，更接近 L4 的先导。 \\
理想的经验是什么 & 数据、算力、组织压力、卫城项目 & 量产自动驾驶是技术、数据、供应链和组织的耦合。 \\
\bottomrule
\end{tabularx}

\begin{warningbox}{时间边界}
视频发布于 2025 年 3 月，但访谈发生在 2024 年 12 月。访谈中说“今年”时多指 2024 年，说“去年”时多指 2023 年。正文按原访谈上下文解释，不把相对时间误读为发布时间。
\end{warningbox}

\subsection{本章小结}

EP96 是张小珺 AI/互联网队列里非常技术化的一集。它可以和 EP120 小鹏 Physical AI、EP132 星海图、VLA 投屏版一起读：自动驾驶不只是汽车行业故事，而是物理世界 AI、数据闭环、端到端模型和世界模型的前哨。

